\documentclass[11pt,a4paper]{article}
\usepackage[utf8]{inputenc}
\usepackage[margin=1in]{geometry}
\input{liatex/preamble}

\usepackage[style=alphabetic]{biblatex}
\addbibresource{notes.bib}

\DeclareMathOperator*{\law}{Law}
\newcommand*{\XX}{\mathcal{X}}
\newcommand*{\RR}{\mathcal{R}}
\newcommand*{\CC}{\mathcal{C}}
\newcommand*{\EE}{\mathcal{E}}
\newcommand*{\lra}{\leftrightarrow}
\DeclareMathOperator*{\Div}{div}

\title{Markov chain notes}
\author{Lily Chung}
\date{}

\begin{document}
\maketitle

Notes on discrete-time, discrete-space, time-homogenous Markov chains, using \cite{LevinPeres} as a reference.

\section*{Definitions}

Let $\XX$ be a countable state space.
A discrete-time stochastic process on $\XX$ is a sequence of random variables $\{X_t : t \in \N\}$ taking values in $\XX$.
Such a process is a \defn{Markov chain} if it satisfies the Markov property, which is that $X_{>t}$ is independent of $X_{<t}$ conditional on $X_t$.
It is \defn{time-homogenous} if there is a Markov kernel $P : S \to S$ such that $\law(X_{t+1}) = P \circ \law(X_t)$ for all $t$.\footnote{Note that $P$ is not necessarily unique if there are states $s$ such that $\Pr[X_t = s] = 0$ for all $t$.}  From now on we consider only time-homogenous Markov chains.  Uses of conditional probability notation should be understood in the context of $P$, so for instance $\Pr[X_1 = y \mid X_0 = x]$ is a synonym for $P[y | x]$ regardless of whether $\Pr[X_0 = x]$ is nonzero.
We will write $\Pr_\pi$ and $\E_\pi$ for probabilities and expectations starting from $X_0 \sim \pi$.

We represent the Markov kernel as a (possibly infinite) row-stochastic matrix $P_{xy} = P[y | x]$,
and we represent distributions on $S$ as row vectors, so that $\pi P$ represents $P \circ \pi$.
Note that functions on $S$ are represented as column vectors, so for instance $\pi P^k f = \E_{\pi}[f(X_k)]$.

For the following definitions, it will be helpful to consider $P$ as a digraph on the vertices $S$, where an edge from $x$ to $y$ exists if $P_{xy} > 0$.

\[
\marginnote{hitting time} \tau_A &\eqdef \inf\{t : X_t \in A\} \\
\marginnote{return time} \tau_A^+ &\eqdef \inf\{t > 0 : X_t \in A\} \\
\marginnote{commute time} \tau_{A,B} &\eqdef \inf\{t > \tau_A : X_t \in B\} \\
t_{x\lra y} &= \E_x \tau_{y,x} = \E_x \tau_y + \E_y \tau_x \\
\marginnote{number of visits before return} v_{AB} &\eqdef \sum_{t = 0}^{\tau_B^+ - 1} [X_t \in A]
\]

\begin{itemize}
\item A Markov chain is \defn{irreducible} if its digraph is strongly connected.
\item A state $x$ is \defn{transient} if $\Pr_x[\tau_x^+ = \infty] > 0$; else \defn{recurrent}.
\item A recurrent state $x$ is \defn{positive recurrent} if $\E_x \tau_x^+ < \infty$; else \defn{null recurrent}.  For finite $S$, all recurrent states are positive recurrent.
\item The \defn{period} of a state $x$ is $\gcd\{n > 0 : \Pr_x[X_n = x] \ne 0\}$.  A state with period 1 is \defn{aperiodic}.
\item A state $i$ is \defn{ergodic} if it is aperiodic and positive recurrent.
\end{itemize}

Transience, positive/null recurrence, period, and ergodicity are constant within an SCC.
For finite $S$ transience is equivalent to existence of an outgoing edge from $i$'s SCC, and null recurrence is impossible.\footnote{Not true for infinite $S$, e.g. a random walk on a 3-dimensional lattice has only a ${\sim}34\%$ chance of ever returning to the origin.}

\section*{Harmonic functions}

A function $f : \XX \to \R$ is \defn{harmonic at $x$} if $(Pf)_x = 0$, and \defn{harmonic} if $Pf = 0$; or equivalently $f(X_t)$ is a martingale.

\begin{itemize}
\item A harmonic, non-negative function on an irreducible recurrent chain must be constant \cite[\nopp 21.4]{LevinPeres}.  
\item Let $B \subset \XX$ be a set of states such that $\tau_B < \infty$ a.s. for all starting states, and let $h_B : B \to \R$ be a bounded function; then $h_B$ extends uniquely to a \emph{bounded} function $h$ which is harmonic at $\XX \setminus B$.  The extension is given by $h(x) = \E_x[h_B(X_{\tau_B})]$.  (Follows from the optional stopping theorem \cite[\nopp 9.1]{LevinPeres}.)
\item Some counterexamples:
  \begin{itemize}
  \item The identity function on $\Z$ is harmonic and non-constant since it takes negative values.
  \item $x \mapsto \E_x[f(X_{\tau_A})]$ on a random walk with absorbing states $A$ is harmonic and non-constant because the chain isn't irreducible.
  \item $x \mapsto \Pr_x[\tau_z^+ < \infty]$ is harmonic, bounded, and non-constant on a transient irreducible chain.
  \item $x, y, z \mapsto (c^z)(-1)^x[x=y]$ where $c + c^{-1} = 6$ is harmonic on $\Z^3$ but vanishes outside the hyperplane $x=y$ \cite[Remark~1.3]{BLMS2017}.
  \end{itemize}
\end{itemize}

\section*{Stationary measures and distributions}

A $\sigma$-finite measure $\mu$ is \defn{stationary} if $\mu P = \mu$.

\begin{itemize}
\item If $z$ is a recurrent state, then $\mu_y = \E_z v_{yz}$ is a stationary measure.
\item An irreducible, recurrent Markov chain has unique stationary measure (up to scaling) given by $\mu_y = \E_z v_{yz}$ \cite[\nopp 21.15]{LevinPeres}.
\end{itemize}

For stationary distributions:

\begin{itemize}
\item A transient or null recurrent state must have zero stationary probability.
\item An irreducible Markov chain has a stationary distribution if and only if it is positive recurrent, in which case the unique stationary distribution is given by $\pi_i = \frac{1}{\E_i \tau_i^+}$ \cite[\nopp 21.12,21.13]{LevinPeres}.
\item In general a distribution over the positive recurrent states is stationary iff each positive recurrent SCC is assigned some multiple of its stationary distribution.
  For finite $S$, the SCCs form a finite DAG,
  the positive recurrent SCCs are the sinks,
  and the stationary distributions are convex combinations of the stationary distributions of each sink.
\item The law of an irreducible ergodic chain converges in TV distance to the unique stationary distribution.  The law of an irreducible null-recurrent or transient chain converges pointwise to the zero measure \cite[\nopp 21.16, 21.19]{LevinPeres}.
\end{itemize}

\subsection*{Time reversal}

\begin{itemize}
\item
  A \defn{time reversal} of $P$ with respect to a measure $\mu$ is a Markov kernel $\hat P$ such that
  \[(\mu P)_x \hat P_{xy} &= \mu_y P_{yx} &\forall x, y.\]
\item $\mu$ \defn{reverses} $P$ (and $P$ is \defn{reversible}) if
  \[\mu_x P_{xy} &= \mu_y P_{yx} &\forall x, y.\]
  This implies $\mu$ is a stationary measure for $P$ and $P$ is its own time reversal with respect to $\mu$.
\end{itemize}

\section*{Electrical networks}

An electrical network is a (possibly-infinite) undirected graph $G$ where each edge $e$ has a \defn{conductance} $c(e) \in [0, \infty)$ such that $c_v < \infty$ for each $v$.

\[
c_v &\eqdef \sum_{u \sim v} c(u, v) \\
c_G &\eqdef \sum_v c_v \\
\marginnote{resistance} r(e) &\eqdef \frac{1}{c(e)} \\
\]

The weighted random walk on an electrical network is a reversible Markov chain: it's reversed by the measure $v \mapsto c_v$, or by the distribution $v \mapsto \frac{c_v}{c_G}$ if $c_G < \infty$.
Conversely a reversible Markov chain is a weighted random walk on the electrical network $c(x, y) = \mu_x P_{xy} = \mu_y P_{yx}$ where $\mu$ is the reversing measure.

Given a set $T \subset \XX$ and a bounded function $V_T : T \to \XX$, the \defn{voltage} is
\[\marginnote{voltage} V(x) &\eqdef \E_x\left[
  \begin{cases}
    V_T(X_{\tau_T}) & \tau_T < \infty \\
    0 & \tau_T = \infty
  \end{cases}
  \right].
\]
By the optional stopping theorem, the voltage is the unique bounded harmonic function which decays to zero at $\infty$.

\subsection*{Current and flows}

For this section only finite networks are considered. \xxx{add a section on $x \lra \infty$ for infinite networks, and generalize the $x \lra y$ section to recurrent infinite networks.}

A \defn{flow} is an antisymmetric real function $\theta$ over the oriented edges, such that $\Div \theta(x) \eqdef \sum_{y \sim x} \theta(x, y) = 0$ for all $x \notin T$.
For a flow $\theta$ from $x$ to $y$ (i.e. $T = \{x, y\}$) define $\norm{\theta} \eqdef \Div \theta(x) = -\Div \theta(y)$.

\[
\marginnote{energy} \EE(\theta) &= \sum_e \theta(e)^2 r(e) \\
\marginnote{current flow} I(x, y) &\eqdef \frac{V(x) - V(y)}{r(x, y)} \\
\marginnote{effective resistance} \RR(x \lra y) &\eqdef \frac{V}{\norm{I}} &\text{on voltage $x \mapsto V, y \mapsto 0$} \\
\marginnote{Thomson's Principle} &= \inf\{\EE(\theta) : \text{$\theta$ a unit flow from $x$ to $y$}\} &\text{uniquely minimized by $\frac{I}{\norm{I}}$} \\
\marginnote{effective conductance} \CC(x \lra y) &\eqdef \frac{1}{\RR(x \lra y)} \\
\marginnote{\cite[\nopp 9.5, 9.6]{LevinPeres}} \Pr_x[\tau_y < \tau_x^+] &= \frac{1}{\E_x v_{yx}} = \frac{1}{c(x) \RR(x \lra y)} & x \ne y \\
\marginnote{\cite[\nopp 10.7]{LevinPeres}} t_{x\lra y} &= c_G \RR(x \lra y)
\]

\section*{Mixing time}

\xxx{todo}

\section*{Cover time}

Assume $S$ is finite.
The \defn{cover time} is the first time when all states in $S$ have been visited.
We are interested in the quantity $C = \max_\pi \E_\pi[\text{cover time}] = \max_v \E_v[\text{cover time}]$.

\begin{itemize}
\item An undirected graph of min-degree $\delta$ has cover time at most $O\left(\frac{nm}{\delta}\right)$ \cite{KLNS1989}.
  Thus for regular graphs it is $O(n^2)$ but it can be $\Theta(n^3)$ for general graphs.
\item An $\Omega(n^3)$ lower bound is the lollipop graph, which is a path of length $\frac{n}{3}$ attached to a clique on $\frac{2n}{3}$ vertices \cite{Feige1995}.
\end{itemize}

\printbibliography

\end{document}
